% =============================================================================
%  Design e arquitetura
%  Cobre a Etapa 4 da trilha: solid, design-of-pattern, ecs, plugin_registry.
% =============================================================================

% =============================================================================
\section{Design e arquitetura}
\label{sec:design}
% =============================================================================

Os capítulos anteriores tratavam de \emph{o que a linguagem faz}. Este trata de
\emph{como organizar o que se escreve com ela} --- e a mudança de assunto vem
com uma mudança de natureza: nada aqui é verificável pelo compilador. Um
programa que viola os cinco princípios SOLID compila, roda e passa nos testes.
O preço é cobrado depois, na primeira mudança de requisito.

A ordem interna é a da trilha, e ela tem uma lógica: \projeto{solid} estabelece
os princípios; \projeto{design-of-pattern} mostra as soluções canônicas que os
realizam; \projeto{ecs} apresenta o contraponto radical à herança; e
\projeto{plugin\_registry} leva o último princípio --- inversão de dependência
--- até o limite do sistema operacional.

% =============================================================================
\subsection{Os cinco princípios SOLID}
\label{sec:solid}
% =============================================================================

O subprojeto \projeto{solid} implementa cada princípio duas vezes: uma versão
que o \emph{viola} e uma que o \emph{cumpre}, lado a lado, executáveis. É a
única forma honesta de discutir o assunto, porque a versão violada quase sempre
é a mais curta --- e o custo dela não aparece até que alguém precise mudar algo.

\subsubsection{SRP --- responsabilidade única}

A definição popular ``uma classe, uma responsabilidade'' é vaga a ponto de ser
inútil. A definição precisa de Robert Martin~\cite{martin2017} usa outra
palavra:

\begin{quote}
\emph{``A module should be responsible to one, and only one, actor.''}
\end{quote}

\textbf{Ator} é um grupo de \emph{stakeholders} que pode exigir a mesma mudança.
Se duas equipes diferentes podem pedir alterações na mesma classe, ela tem duas
responsabilidades --- e a mudança pedida por uma vai quebrar o que a outra
esperava.

O exemplo canônico do livro, reproduzido no subprojeto, é uma classe
\texttt{Employee} que serve a três atores:

\begin{lstlisting}[style=bt, language=bash,
    caption={Saída de \texttt{src/solid/srp/}: o acoplamento oculto.}]
Employee serve a 3 atores:
   calculatePay()  → COO / RH
   reportHours()   → CFO (auditoria)
   save()          → CTO / DBA

   Risco: total_hours() compartilhado entre calculatePay() e reportHours()
   Uma mudança pedida pelo COO afeta silenciosamente o relatório do CFO!
\end{lstlisting}

O perigo não é a classe ser grande. É o método privado \texttt{total\_hours()}
ser \emph{compartilhado} entre dois atores que nunca conversaram: quando o COO
pedir uma mudança na regra de horas, o relatório de auditoria do CFO muda
junto, e ninguém será avisado.

\begin{regra}[a heurística do ``e'']
Se você precisa da conjunção ``e'' ou ``ou'' para descrever o que uma função
faz, ela tem mais de uma responsabilidade. \texttt{process\_and\_print\_student}
é o nome delatando o problema. O complemento natural para funções é a
\emph{Command-Query Separation}: um \emph{command} muda estado e devolve
\texttt{void}; uma \emph{query} lê estado e não tem efeito colateral. Um
\texttt{int increment\_and\_get()} é os dois ao mesmo tempo.
\end{regra}

\begin{table}[htbp]
\centering
\caption{Sinais de violação de SRP, e o que cada um costuma indicar.}
\label{tab:srp-sinais}
\footnotesize
\begin{tabular}{@{}ll@{}}
\toprule
\textbf{Sinal} & \textbf{Problema provável} \\
\midrule
classe com mais de 200--300 linhas                       & múltiplas responsabilidades \\
nome genérico: \texttt{Manager}, \texttt{Helper}, \texttt{Utils} & \emph{god class} \\
método com ``e'' ou ``ou'' na descrição                  & duas responsabilidades \\
teste que precisa de 5 ou mais \emph{mocks}              & a classe faz demais \\
mudar formato de e-mail quebra cálculo de salário        & acoplamento oculto \\
classe inclui \texttt{<filesystem>}, \texttt{<regex>} \emph{e} \texttt{<curl.h>} & mistura de camadas \\
\bottomrule
\end{tabular}
\end{table}

\begin{armadilha}[SRP também pode ser aplicado demais]
Classes de uma linha, um \emph{getter} por classe, módulos tão pequenos que o
custo de coordená-los supera o benefício --- todos são sintomas de
sobre-aplicação. A regra de ouro é \textbf{coesão acima de granularidade}:
agrupe o que muda junto, separe o que muda por razões diferentes.
\end{armadilha}

\subsubsection{OCP --- aberto para extensão, fechado para modificação}

O sintoma de violação é sempre o mesmo: um \texttt{switch} ou uma cascata de
\texttt{if/else if} sobre um \emph{tipo}, repetida em várias funções.

\begin{lstlisting}[style=bt, language=C++,
    caption={A cascata, e o motivo de ela não escalar.}]
double area(const Forma& f) {
    if      (f.tipo == Circulo)   return 3.14159 * f.raio * f.raio;
    else if (f.tipo == Retangulo) return f.largura * f.altura;
    // ← e um `else if` por forma nova. Para sempre.
}
double perimetro(const Forma& f) { /* a MESMA cascata, repetida */ }
\end{lstlisting}

Acrescentar um triângulo exige editar \textbf{quatro} lugares já testados --- e
esquecer um deles não gera erro de compilação: gera resultado errado. Com
polimorfismo, o triângulo é um \emph{arquivo novo}, a função que consome nunca é
tocada, e o \texttt{= 0} obriga a implementar tudo.

O subprojeto conta o custo, e a linha que decide é a terceira:

\begin{table}[htbp]
\centering
\caption{O que muda ao acrescentar um caso novo. Código que não é tocado não
         pode regredir.}
\label{tab:ocp-custo}
\footnotesize
\begin{tabular}{@{}lcc@{}}
\toprule
\textbf{Ao acrescentar um caso} & \textbf{Com \texttt{switch}} & \textbf{Com OCP} \\
\midrule
arquivos criados          & 0      & 1 \\
arquivos \textbf{editados}& 3--5   & \textbf{0} \\
testes a reescrever       & 3--5   & \textbf{0} \\
risco de regressão        & alto   & nenhum \\
\bottomrule
\end{tabular}
\end{table}

Quando o que varia é um \emph{algoritmo} --- e não uma família de entidades ---
a forma correta é injetá-lo, e há quatro mecanismos com trocas diferentes
(\cref{tab:ocp-mecanismos}).

\begin{table}[htbp]
\centering
\caption{Quatro formas de injetar variação. ``Homogêneo'' significa poder
         guardar variantes no mesmo \texttt{vector}.}
\label{tab:ocp-mecanismos}
\footnotesize
\begin{tabular}{@{}llcc@{}}
\toprule
\textbf{Mecanismo} & \textbf{Extensão} & \textbf{Custo} & \textbf{Homogêneo?} \\
\midrule
herança                  & execução   & \gls{vtable}   & sim \\
\texttt{std::function}   & execução   & indireção      & sim \\
\emph{template} (policy) & compilação & \textbf{zero}  & não \\
\emph{concept}           & compilação & \textbf{zero}  & não \\
\bottomrule
\end{tabular}
\end{table}

O \emph{concept} resolve o pior defeito da \emph{policy} por \emph{template}, que
é a mensagem de erro. Declarando o contrato explicitamente, uma política
incompleta falha \emph{na linha da chamada}, dizendo qual expressão faltou ---
em vez de trezentas linhas de instanciação (\cref{sec:concepts}):

\begin{lstlisting}[style=bt, language=C++,
    caption={O contrato da política, escrito como \emph{concept}.}]
template <typename P>
concept PoliticaDeDesconto = requires(double total) {
    { P::aplicar(total) } -> std::convertible_to<double>;
    { P::nome() }         -> std::convertible_to<std::string>;
};

template <PoliticaDeDesconto P> class Carrinho { /* ... */ };
\end{lstlisting}

O último demo do subprojeto vai além: um \textbf{registro auto-populado} elimina
até a edição do ponto de composição. Uma macro no arquivo do exportador basta
--- \texttt{REGISTRAR\_EXPORTADOR(ExportadorMarkdown, "markdown")} --- e o laço
que consome nunca foi editado, nem conhece os nomes das classes.

\begin{armadilha}[o registro auto-populado tem três armadilhas de ligação]
Primeira: a ordem de inicialização estática entre unidades de tradução não é
definida, então o registro precisa ser um \emph{Meyers singleton} --- função com
\texttt{static} local --- e nunca uma variável global. Segunda: o \emph{linker}
pode \textbf{descartar} o objeto estático se nada naquela unidade for
referenciado (\texttt{-{}-gc-sections}, bibliotecas estáticas), o que exige
\texttt{-{}-whole-archive} ou um símbolo âncora. Terceira: o fluxo fica indireto
--- não se descobre quem registrou lendo o \texttt{main}.
\end{armadilha}

\begin{regra}[Rule of Three: OCP responde ao passado, não profetiza o futuro]
OCP aplicado \textbf{antes} de existir um segundo caso é custo puro --- e a
abstração escolhida com uma amostra de tamanho um quase sempre é a errada. Na
primeira vez, escreva direto. Na segunda, duplique e anote o incômodo. Na
terceira, abstraia --- porque agora você já viu o \emph{eixo} da variação.
E há o caso em que a cascata é a resposta certa: quando o conjunto é
\emph{fechado} (quatro naipes, três estados de semáforo), o \texttt{switch} é
mais simples e mais rápido, ou então \texttt{std::variant} com \texttt{visit}
(\cref{sec:vocabulary}). A pergunta que decide: \emph{casos novos vão chegar
depois do \emph{deploy}, vindos de quem não controla este arquivo?}
\end{regra}

\subsubsection{LSP --- substituição de Liskov}

É o terceiro princípio e o único que o compilador \textbf{não pode verificar}.
LSP é sobre \emph{comportamento}, não sobre assinatura: o compilador confere a
assinatura, e ninguém confere o comportamento. Por isso violações de LSP
compilam perfeitamente e aparecem em produção.

O exemplo canônico é canônico porque ``quadrado é um retângulo'' é verdade na
matemática e mentira no código:

\begin{lstlisting}[style=bt, language=bash,
    caption={Saída de \texttt{src/solid/lsp/}: a mesma sequência de chamadas,
             resultados diferentes.}]
com Retangulo: largura=5, altura=4 → área esperada 20, obtida 20  ✓
com Quadrado:  largura=5, altura=4 → área esperada 20, obtida 16  ✗ LSP VIOLADO
\end{lstlisting}

O erro não foi implementar mal o quadrado. Foi afirmar que ele \emph{é um}
retângulo \textbf{mutável} --- e essa relação não sobrevive à mutabilidade. As
duas saídas boas são a imutabilidade (sem \emph{setters}, o conflito desaparece)
e tratá-los como irmãos sob uma interface que ambos \emph{podem} cumprir
(\texttt{escala()} em vez de \emph{setters} independentes).

\begin{regra}[a pergunta certa não é ``X é um Y?'']
É: \emph{X consegue cumprir \textbf{todas} as promessas que Y faz aos seus
clientes?} A primeira pergunta é sobre taxonomia; a segunda é sobre contrato ---
e só a segunda tem consequência para o código.
\end{regra}

O contrato se decompõe em quatro regras, e todas seguem a mesma direção
(\cref{tab:lsp-regras}).

\begin{table}[htbp]
\centering
\caption{O derivado pode ser mais \emph{generoso}, nunca mais \emph{exigente} ---
         a mesma contravariância dos tipos de função.}
\label{tab:lsp-regras}
\footnotesize
\begin{tabular}{@{}lccc@{}}
\toprule
\textbf{O subtipo pode\ldots} & \textbf{Pré-condição} & \textbf{Pós-condição} & \textbf{Invariante} \\
\midrule
exigir \textbf{mais}    & não & --- & --- \\
exigir \textbf{menos}   & sim & --- & --- \\
entregar \textbf{mais}  & --- & sim & sim \\
entregar \textbf{menos} & --- & não & não \\
\bottomrule
\end{tabular}
\end{table}

As três primeiras regras são diretas: fortalecer uma pré-condição quebra quem
chamava com argumentos que a base aceitava (um depósito ``VIP'' que exige mínimo
de 100); enfraquecer uma pós-condição quebra quem contava com ela (um
``ordenador'' que também deduplica devolve menos elementos do que entraram);
violar uma invariante sabota todo código que confiava nela (uma conta com cheque
especial cujo saldo fica negativo).

\begin{armadilha}[a \emph{history rule} não muda nenhuma assinatura]
A quarta regra é a menos conhecida: um derivado não pode permitir mudanças de
estado que a base proibia.
\begin{lstlisting}[style=bt, language=bash, frame=none,
                   xleftmargin=0pt, xrightmargin=0pt]
PontoImutavel  → cache válido ✓
PontoMutavel   → cache INVÁLIDO ✗ LSP
\end{lstlisting}
Nenhum método foi sobrescrito, nenhuma assinatura mudou. O que mudou foi o
conjunto de \textbf{histórias possíveis} do objeto --- e quem tinha um cache
baseado na imutabilidade passou a ter um cache errado. É por isso que
\texttt{final} e a imutabilidade são ferramentas de LSP, e não apenas de
otimização.
\end{armadilha}

Como o compilador não ajuda, o subprojeto oferece dois instrumentos práticos. O
primeiro é o \textbf{teste de contrato}: uma suíte escrita contra a \emph{base} e
rodada contra \emph{cada} derivado --- em GoogleTest, é o que
\texttt{TYPED\_TEST\_SUITE} faz.

\begin{lstlisting}[style=bt, language=bash,
    caption={Saída de \texttt{src/solid/lsp/}: o LSP virando código executável.}]
PilhaVetor         5 passou, 0 falhou  ✓
PilhaLimitada(3)   3 passou, 2 falhou  ✗ LSP VIOLADO
  ✗ empilhar 10 → tamanho 10
  ✗ round-trip: 5 entraram, 5 devem sair
\end{lstlisting}

O segundo é o idioma \textbf{NVI}, já visto na \cref{sec:oop} como solução para
argumentos padrão em virtuais --- e que aqui aparece no seu papel principal:
colocar o contrato na base, onde o derivado não pode afrouxá-lo.

\begin{lstlisting}[style=bt, language=C++,
    caption={NVI: pré e pós-condições na fronteira pública, não virtual.}]
std::string buscar(int id) {              // público, NÃO virtual
    if (id <= 0) throw std::invalid_argument("id deve ser positivo");  // pré
    std::string r = buscar_impl(id);                                    // miolo
    if (r.empty()) throw std::runtime_error("devolveu vazio");          // pós
    return r;
}
private:
    virtual std::string buscar_impl(int id) = 0;   // privado, virtual
\end{lstlisting}

\subsubsection{ISP --- segregação de interfaces}

ISP e LSP são o mesmo problema visto de dois lados: LSP pergunta \emph{o derivado
consegue cumprir tudo que a base promete?}; ISP pergunta \emph{a base está
prometendo coisas demais?}

A \emph{fat interface} força um implementador a declarar métodos que ele não sabe
fazer, e as três saídas disponíveis são todas ruins
(\cref{tab:isp-saidas}).

\begin{table}[htbp]
\centering
\caption{O que fazer num \texttt{override} que não se sabe implementar --- e por
         que nenhuma das opções serve.}
\label{tab:isp-saidas}
\footnotesize
\begin{tabular}{@{}ll@{}}
\toprule
\textbf{Saída} & \textbf{Problema} \\
\midrule
lançar exceção          & viola o LSP --- o cliente quebra \\
não fazer nada          & pior: quebra em \textbf{silêncio} \\
devolver valor inventado& pior ainda: corrompe dados \\
\bottomrule
\end{tabular}
\end{table}

A correção segrega cada capacidade em uma interface própria, e o ganho decisivo
é de categoria: \textbf{o que era exceção em tempo de execução (ou silêncio)
vira erro de compilação}.

\begin{lstlisting}[style=bt, language=C++,
    caption={O cliente pede o mínimo de que precisa.}]
void fluxo_de_impressao(Imprimivel& d, const std::string& doc);
void fluxo_de_arquivamento(Digitalizavel& d, const std::string& destino);

fluxo_de_impressao(scanner, "x");    // ERRO DE COMPILAÇÃO
\end{lstlisting}

\begin{porque}[a interface pertence ao cliente, não ao fornecedor]
A leitura ingênua do ISP é ``interfaces pequenas''. A leitura correta é mais
forte: a interface é definida pelo que o \textbf{cliente} precisa, e por isso
deve ser \emph{declarada junto do cliente}. Um serviço de doze métodos pode ser
consumido por dez clientes através de dez interfaces minúsculas, cada uma
nomeada segundo o \emph{papel} que cumpre, com adaptadores fininhos ligando o
serviço legado a cada papel. O que isso compra: o enviador de e-mails
\textbf{não consegue} tocar em dinheiro, garantido pelo tipo; e os dublês de
teste têm dois métodos cada, em vez de doze --- sem \emph{framework} de
\emph{mock}, sem banco, sem rede.
\end{porque}

Daí decorre uma regra de nomenclatura que muda o desenho: nomeie pelo
\textbf{papel}, não pela classe concreta. \texttt{FonteDeContato} em vez de
\texttt{IServicoDeUsuarios}; \texttt{LeitorDePedidos} em vez de
\texttt{IRepositorioMySql}; qualquer coisa em vez de \texttt{IHelper},
\texttt{IManager} ou \texttt{IUtil}.

\begin{armadilha}[segregar demais vira sopa de interfaces]
O limite é objetivo: cada interface deve ter \emph{pelo menos um cliente que use
todos os seus métodos}. Se nenhum cliente usa a interface inteira, ela ainda não
foi segregada o bastante; se cada interface tem um método e nenhum cliente
real, foi segregada demais.
\end{armadilha}

\subsubsection{DIP --- inversão de dependência}

O nome engana. \textbf{``Inversão'' não é sobre injetar dependência} --- é sobre
a direção da seta no diagrama de módulos (\cref{fig:dip}).

\begin{figure}[htbp]
\centering
\begin{tikzpicture}[node distance=6mm and 14mm]
  \node[caixarisco] (pol1) {Política de negócio};
  \node[caixarisco, right=of pol1] (db1) {MySQL};
  \draw[setarisco] (pol1) -- node[rotulo,above] {depende} (db1);

  \node[caixaok, below=14mm of pol1] (pol2) {Política de negócio};
  \node[folha, right=of pol2] (abs) {\texttt{IRepositorio}\\\scriptsize declarada na política};
  \node[caixaok, right=of abs] (db2) {MySQL};
  \draw[setaok] (pol2) -- (abs);
  \draw[setaok] (db2) -- node[rotulo,above] {implementa} (abs);

  \node[rotulo, left=2mm of pol1] {antes};
  \node[rotulo, left=2mm of pol2] {depois};
\end{tikzpicture}
\caption{A inversão. A seta do MySQL passou a apontar \emph{para dentro}, na
         direção da abstração que a própria política definiu --- e é a
         abstração morar na camada de negócio, e não numa camada neutra, que
         caracteriza a inversão.}
\label{fig:dip}
\end{figure}

O detalhe que faz a diferença é onde a abstração é \emph{declarada} e como é
\emph{nomeada}: \texttt{RepositorioDePedidos}, ao lado da política, e não
\texttt{IBancoDeDados}, ao lado da infraestrutura. O mesmo serviço passa a rodar
com MySQL, com arquivo e em memória --- e nunca soube que MySQL existe.

\begin{regra}[o teste decisivo é a direção do \texttt{\#include}]
Em que direção o código-fonte aponta em tempo de compilação? Se o módulo de
negócio faz \texttt{\#include "mysql.hpp"}, não houve inversão nenhuma --- por
mais interfaces que existam no meio. A verificação é mecânica:
\texttt{grep -r '\#include' src/entidades/}. Se aparecer um cabeçalho de SQL, a
arquitetura foi violada.
\end{regra}

As formas de injeção diferem pelo momento em que a dependência é fixada:
construtor para dependência obrigatória (o estado inválido fica inexprimível);
parâmetro para o que varia a cada chamada; \emph{template} quando a escolha é
conhecida em compilação e o custo deve ser zero. Nesta última, o \emph{concept}
\textbf{é} a abstração:

\begin{lstlisting}[style=bt, language=C++,
    caption={Injeção estática: a dependência continua invertida, e não custa nada.}]
template <RelogioEstatico R>
class AuditoriaEstatica {
    R relogio_;                       // por VALOR: sem indireção, sem heap
};
// sizeof(AuditoriaEstatica<RelogioRapido>) = 1 byte
\end{lstlisting}

Se o serviço não pode construir o repositório, alguém precisa construir --- e
esse alguém é o \textbf{composition root}: um ponto único que pode conhecer tudo
porque não tem regra de negócio nenhuma. O subprojeto o materializa em duas
funções, \texttt{montar\_producao()} e \texttt{montar\_teste()}, e a
consequência é que \emph{todo \texttt{new} do programa está nessas duas
funções}.

\begin{armadilha}[\emph{service locator} não é injeção --- é pior que o acoplamento direto]
\begin{lstlisting}[style=bt, language=C++, frame=none,
                   xleftmargin=0pt, xrightmargin=0pt]
void fazer() {
    auto* log = Locator::obter<Log>();   // ← o objeto vai BUSCAR
}
\end{lstlisting}
As dependências ficam \textbf{invisíveis na assinatura}, o teste vira estado
global mutável, e a falha é em tempo de execução --- ao passo que a injeção por
construtor falha em compilação. A regra que separa os dois: na injeção, a
dependência é \textbf{empurrada} para o objeto; no \emph{locator}, o objeto vai
\textbf{buscar}. Se a classe busca, ela conhece o buscador, e não houve
inversão.
\end{armadilha}

Em escala, DIP vira a \emph{regra da dependência}: o código-fonte só pode
apontar para dentro, na direção das políticas de mais alto nível
(\cref{tab:camadas}). E o ponto que dá nome ao princípio é que o fluxo de
\emph{controle} vai de dentro para fora --- o caso de uso chama o banco ---
enquanto o fluxo de \emph{código} vai de fora para dentro. São opostos.

\begin{table}[htbp]
\centering
\caption{A regra da dependência: cada círculo só inclui os de dentro.}
\label{tab:camadas}
\footnotesize
\begin{tabular}{@{}lll@{}}
\toprule
\textbf{Círculo} & \textbf{Pode incluir} & \textbf{Volatilidade} \\
\midrule
entidades    & \textbf{nada}                & mínima \\
casos de uso & entidades                    & baixa \\
adaptadores  & casos de uso + entidades     & média \\
\emph{frameworks} & tudo                    & máxima \\
\bottomrule
\end{tabular}
\end{table}

\begin{regra}[inverta o volátil; não inverta o estável]
DIP paga-se com uma interface, uma indireção e um arquivo por dependência. Não
vale quando a dependência é \emph{estável} --- ninguém abstrai
\texttt{std::string}, \texttt{std::vector} ou \texttt{fmt} --- nem quando o
programa é pequeno e sem testes automatizados, nem quando existe um
implementador e não haverá outro. A pergunta que decide: \emph{esta dependência
é volátil?} Inverta banco, rede, relógio, sistema de arquivos e aleatoriedade.
Não inverta o que é estável há trinta anos.
\end{regra}

% =============================================================================
\subsection{Padrões de projeto}
\label{sec:patterns}
% =============================================================================

O subprojeto \projeto{design-of-pattern} implementa \textbf{24 padrões}: os 23
do catálogo GoF~\cite{gof1994} e mais o \emph{Null Object}, que não está no
livro e aparece em quase todo código real. Cada padrão mora no seu diretório,
tem \texttt{meson.build} próprio e produz um executável --- são 24 alvos, todos
construídos e todos executáveis a partir de \texttt{dist/bin/}. Nenhum está
comentado, e os números citados adiante foram lidos da saída desses binários,
não de livro.

Antes de tudo, uma definição, porque a palavra costuma ser usada de forma
descuidada. Um \textbf{padrão de projeto} não é uma biblioteca, não é uma
classe-base para herdar e não é uma receita a seguir: é o \emph{nome} de uma
solução que foi redescoberta de forma independente tantas vezes que valeu a
pena catalogá-la. O valor está tanto na solução quanto no nome --- dizer ``isto
é um \emph{Strategy}'' comunica em duas palavras o que levaria um parágrafo, e
é por isso que o catálogo sobreviveu a três revisões da linguagem.

\begin{porque}[por que estudar 24 padrões se metade deles ``sumiu'' no C++ moderno?]
Porque o que ``sumiu'' foi a \emph{implementação}, não o \emph{problema}. O
\emph{Strategy} continua existindo quando um \classe{std::function} substitui a
hierarquia --- só que agora o padrão tem uma linha em vez de quatro classes. E
saber que aquilo é um \emph{Strategy} é o que permite reconhecer, seis meses
depois, que a variação é do algoritmo e não do tipo. Um padrão é vocabulário
antes de ser código: o dia em que o \classe{std::function} não servir --- porque
o algoritmo precisa de estado, de nome, de ciclo de vida --- é o padrão clássico
que volta, e ele precisa estar aprendido.
\end{porque}

\begin{table}[htbp]
\centering
\caption{Os 24 padrões implementados em \texttt{src/design-of-pattern/}, e o que
         o executável de cada um demonstra. A coluna da direita reproduz o que o
         \texttt{README.md} do subprojeto registra como objetivo de cada demo.}
\label{tab:padroes}
\footnotesize
\begin{tabular}{@{}llp{0.44\textwidth}@{}}
\toprule
\textbf{Padrão} & \textbf{Categoria} & \textbf{O que o demo mostra} \\
\midrule
\emph{Factory Method}  & criacional & criação polimórfica por nome \\
\emph{Abstract Factory}& criacional & famílias compatíveis; misturar temas fica inexprimível \\
\emph{Builder}         & criacional & \texttt{construir() \&\&} impede reusar o \emph{builder} \\
\emph{Prototype}       & criacional & clone profundo medido; CRTP elimina o \emph{boilerplate} \\
\emph{Singleton}       & criacional & \emph{Meyers singleton} e o custo dele \\
\midrule
\emph{Adapter}   & estrutural & API C legada $\rightarrow$ \texttt{expected} + \gls{raii} \\
\emph{Bridge}    & estrutural & $3\times3 = 9$ classes vira $3+3$; \emph{pimpl} é um \emph{Bridge} degenerado \\
\emph{Composite} & estrutural & \texttt{adicionar} só no composto, por LSP/ISP \\
\emph{Decorator} & estrutural & empilhamento de comportamento \\
\emph{Facade}    & estrutural & 6 subsistemas $\rightarrow$ 1 chamada; o subsistema segue acessível \\
\emph{Flyweight} & estrutural & \textbf{10,4$\times$ menos memória}, medido \\
\emph{Proxy}     & estrutural & virtual/proteção/\emph{smart}; \classe{std::shared\_ptr} é um deles \\
\midrule
\emph{Chain of Responsibility} & comportamental & \emph{pipeline} de \emph{handlers} \\
\emph{Command}         & comportamental & \emph{undo} por operação inversa \\
\emph{Interpreter}     & comportamental & AST com hierarquia virtual \\
\emph{Iterator}        & comportamental & iteradores próprios, direto e reverso \\
\emph{Mediator}        & comportamental & teia $N\times(N-1)$ vira estrela $N$ \\
\emph{Memento}         & comportamental & opacidade real, via \texttt{friend} \\
\emph{Observer}        & comportamental & notificação $1\rightarrow N$ \\
\emph{State}           & comportamental & máquina de estados de TCP \\
\emph{Strategy}        & comportamental & algoritmos intercambiáveis \\
\emph{Template Method} & comportamental & obrigatório \emph{vs.} \emph{hook} \emph{vs.} fixo \\
\emph{Visitor}         & comportamental & despacho duplo, e o \texttt{variant} fazendo melhor \\
\midrule
\emph{Null Object}     & fora do catálogo & eliminar verificações de \texttt{nullptr} espalhadas \\
\bottomrule
\end{tabular}
\end{table}

% -----------------------------------------------------------------------------
\subsubsection{O eixo que decide quase tudo: o \emph{expression problem}}
\label{sec:expression-problem}
% -----------------------------------------------------------------------------

Vários padrões do catálogo são respostas à \emph{mesma} pergunta, formulada de
maneiras diferentes. Reconhecer a pergunta economiza metade do estudo.

Considere um programa que manipula uma coleção de \textbf{tipos} (nós de uma
árvore sintática, formas geométricas, widgets) sobre os quais se executam
\textbf{operações} (imprimir, avaliar, serializar, calcular área). Há duas
direções em que esse programa pode crescer, e o \emph{expression
problem} é a observação de que \textbf{as duas ferramentas
disponíveis são boas em direções opostas}.

\begin{porque}[o que exatamente é ``aberto'' e ``fechado'' aqui]
Com uma \textbf{hierarquia virtual}, acrescentar um \emph{tipo} é criar um
arquivo: a classe nova herda da base, implementa os métodos, e nenhum código
existente é tocado. Mas acrescentar uma \emph{operação} é editar a base e
\emph{todas} as derivadas --- e recompilar tudo.

Com um \classe{std::variant}, é o inverso exato: acrescentar uma \emph{operação}
é escrever uma função livre nova, sem tocar em nada; mas acrescentar um
\emph{tipo} à lista do \texttt{variant} obriga a revisitar todos os
\texttt{visit} do programa. A diferença decisiva é que, no segundo caso,
\textbf{o compilador aponta cada lugar} --- um \texttt{visit} que não cobre
todas as alternativas não compila. Na hierarquia, esquecer de sobrescrever um
método novo com corpo padrão compila e falha em execução.
\end{porque}

\begin{table}[htbp]
\centering
\caption{O \emph{expression problem}, como o \texttt{README.md} de
         \texttt{design-of-pattern} o formula: qual eixo do seu problema cresce?}
\label{tab:expression-problem}
\footnotesize
\begin{tabular}{@{}lll@{}}
\toprule
\textbf{Cresce mais} & \textbf{Use} & \textbf{Padrões que realizam} \\
\midrule
os \textbf{tipos}     & hierarquia virtual        & \emph{Visitor}, \emph{Composite}, \emph{Abstract Factory} \\
as \textbf{operações} & \texttt{variant} + \texttt{visit} & --- (função livre, exaustividade verificada) \\
\bottomrule
\end{tabular}
\end{table}

\begin{regra}[a pergunta que substitui metade da discussão sobre padrões]
\emph{O que vai variar mais no meu problema: os tipos ou as operações?} Se a
resposta for ``os tipos, e eles vêm de fora --- de \emph{plugins}, de
bibliotecas que não controlo'', a hierarquia virtual é a resposta certa e
continua sendo. Se for ``os tipos são poucos e conhecidos, mas toda semana
alguém pede uma operação nova'', o \classe{std::variant} é a resposta certa e
ainda dá exaustividade de graça. Nenhum dos dois é melhor em abstrato.
\end{regra}

% -----------------------------------------------------------------------------
\subsubsection{Os cinco criacionais}
% -----------------------------------------------------------------------------

Os padrões criacionais respondem a uma pergunta só: \emph{quem decide qual
classe concreta será instanciada, e quando?} A resposta ingênua --- o cliente,
com \texttt{new} --- acopla o cliente a todos os tipos concretos e é o que os
cinco evitam, cada um por um motivo diferente.

\paragraph{\emph{Factory Method}: criar sem nomear a classe.} O cliente pede
\texttt{ShapeFactory::create("circle", \{5.0\})} e recebe um
\texttt{unique\_ptr<Shape>}. O \texttt{switch} sobre a \emph{string} existe, mas
existe em \textbf{um} lugar --- e esse lugar é o único do programa que conhece
os nomes concretos. É a mesma ideia do \emph{composition root} da
\cref{sec:solid}, aplicada a uma família só.

\paragraph{\emph{Abstract Factory}: a consistência da família virando tipo.} A
diferença em relação ao anterior costuma escapar, e o demo a torna visível numa
linha:

\begin{lstlisting}[style=bt, language=C++,
    caption={Resumo de \texttt{src/design-of-pattern/abstract\_factory/main.cpp}: o que a
             \emph{Factory} simples não impede.}]
// Com Factory simples, uma por produto:
auto b = criar_botao("escuro");
auto c = criar_campo("claro");    // COMPILA -- e a UI fica inconsistente

// Com Abstract Factory:
const FabricaDeWidgets& f = escolher_tema();
auto b = f.criar_botao("Salvar");
auto c = f.criar_campo();         // MESMA familia, garantido pelo TIPO
\end{lstlisting}

A família é escolhida uma vez; depois disso, misturar temas não é
``desencorajado'' --- é \textbf{inexprimível}. E a função
\texttt{renderizar\_formulario} foi escrita uma vez, contra
\texttt{FabricaDeWidgets}: acrescentar um terceiro tema não muda uma linha dela.

\paragraph{\emph{Builder}: construção com validação no fim.} O problema aparece
antes da solução. Sem \emph{Builder}, há três saídas, e as três são ruins:

\begin{lstlisting}[style=bt, language=bash,
    caption={Saída de \texttt{dist/bin/builder}, seção 1: as três saídas e seus
             defeitos.}]
(a) construtor gigante:
      Relatorio r{"Vendas","csv",true,false,true,100,",","utf-8",2};
    -> trocar dois `bool` de lugar COMPILA e muda o comportamento

(b) N sobrecargas:
      Relatorio(t, f);  Relatorio(t, f, cabecalho);  ...
    -> 2^n combinacoes para n campos opcionais

(c) setters publicos:
      Relatorio r; r.set_titulo(...); r.set_formato(...);
    -> entre as chamadas, o objeto esta INVALIDO e isso e visivel
\end{lstlisting}

O \emph{Builder} resolve as três: os campos obrigatórios estão no construtor do
\emph{builder} (não há como esquecê-los, porque o \emph{builder} não existe sem
eles), a ordem dos opcionais é livre, e cada passo tem \emph{nome} ---
\texttt{.com\_rodape()} é mais legível que um \texttt{true} na quarta posição.
E a validação acontece uma vez, em \texttt{construir()}, quando todos os campos
já foram definidos:

\begin{lstlisting}[style=bt, language=bash,
    caption={Saída de \texttt{dist/bin/builder}, seção 3: regras que dependem de
             \emph{dois} campos.}]
formato valido       -> ok: 'X' [html] cab=sim rod=nao sep=',' dec=2
formato invalido     -> recusado: formato deve ser csv, json ou html
casas_decimais = 99  -> recusado: casas_decimais fora de 0..10
csv sem separador    -> recusado: csv exige um separador
\end{lstlisting}

A última linha é a que justifica o padrão: um \emph{setter} público jamais
conseguiria validar ``csv exige separador'', porque no momento em que ele roda
não se sabe se o outro campo já foi preenchido.

\begin{armadilha}[o \texttt{\&\&} nos métodos do \emph{builder} não é enfeite]
Todos os métodos do \emph{builder} do subprojeto são qualificados com
\texttt{\&\&} --- inclusive \texttt{construir()}:
\begin{lstlisting}[style=bt, language=C++, frame=none,
                   xleftmargin=0pt, xrightmargin=0pt]
Construtor&& limite(int n) &&  { ...; return std::move(*this); }
Relatorio    construir()   &&;
\end{lstlisting}
O qualificador de referência (\cref{sec:value-semantics}) faz com que os
métodos só possam ser chamados sobre um \emph{rvalue}. A consequência prática é
que o \emph{builder} \textbf{não pode ser reusado depois de construir}: o objeto
foi movido para dentro do produto, e chamar outro método nele seria ler um
objeto ``movido de''. O compilador passa a recusar exatamente o erro que a
versão sem \texttt{\&\&} permitiria silenciosamente.
\end{armadilha}

\paragraph{\emph{Prototype}: copiar sem conhecer o tipo concreto.} A assinatura
é \texttt{virtual std::unique\_ptr<Forma> clonar() const}, e o detalhe que o
demo mede é que a cópia precisa ser \textbf{profunda}: um nó composto que copie
apenas os ponteiros produz duas árvores compartilhando filhos. No caso do
subprojeto, o código nem compilaria --- \classe{std::unique\_ptr} não é
copiável, e é o próprio sistema de tipos denunciando a cópia rasa.

Como todo \texttt{clonar()} é literalmente a mesma linha, uma camada CRTP
(\cref{sec:templates}) escreve por todas:

\begin{lstlisting}[style=bt, language=C++,
    caption={De \texttt{src/design-of-pattern/prototype/}: o CRTP eliminando o
             \emph{boilerplate}. Nenhuma das classes concretas escreveu
             \texttt{clonar()}.}]
template <typename Derivada, typename Base = Forma>
class Clonavel : public Base {
    std::unique_ptr<Forma> clonar() const override {
        return std::make_unique<Derivada>(static_cast<const Derivada&>(*this));
    }
};
class Quadrado : public Clonavel<Quadrado> { /* ... */ };
\end{lstlisting}

\paragraph{\emph{Singleton}.} Discutido adiante (\cref{sec:padroes-modernos}),
porque nele a ``versão moderna'' resolveu um problema e deixou o outro intacto.

\begin{regra}[o \emph{Prototype} ganha da \emph{Factory} num caso específico]
Uma \emph{factory} só sabe criar o que foi \textbf{compilado} nela. Um registro
de protótipos guarda \emph{instâncias-modelo} que podem ter sido configuradas em
tempo de execução --- carregadas de um arquivo, montadas pelo usuário --- e
clonadas sob demanda. É esse uso que dá nome ao padrão, e é a única situação em
que ele vence claramente a \emph{factory}.
\end{regra}

% -----------------------------------------------------------------------------
\subsubsection{Os sete estruturais}
% -----------------------------------------------------------------------------

Os estruturais respondem a \emph{como compor objetos em estruturas maiores sem
que o cliente precise saber da composição}. Quatro deles --- \emph{Adapter},
\emph{Decorator}, \emph{Proxy} e \emph{Facade} --- têm praticamente o mesmo
diagrama de classes, e a diferença entre eles é de \emph{intenção}; a
\cref{tab:confusoes} adiante separa os quatro de uma vez.

\paragraph{\emph{Adapter}: quando a interface é a errada.} O demo adapta uma
biblioteca C legada, que não pode ser modificada, a uma interface moderna. São
quatro incompatibilidades simultâneas, e o adaptador resolve as quatro:

\begin{lstlisting}[style=bt, language=bash,
    caption={Saída de \texttt{dist/bin/adapter}, seção 1: o que o adaptador
             traduz.}]
codigo de erro int    -> std::expected
int magico de nivel   -> enum class Nivel
(char*, size_t)       -> std::string_view
create/destroy manual -> RAII
\end{lstlisting}

O ganho fica visível na terceira seção do demo: o \texttt{LEGACY\_BUFFER\_FULL}
--- um \texttt{int} de valor 2 --- vira \texttt{std::unexpected\{ErroDeLog::
BufferCheio\}}, e o cliente trata com \texttt{if (r)} sem nunca ver o código
numérico (\cref{sec:errors}). A função de negócio recebe \texttt{adapter::Log\&}
e nunca vai saber que existe uma biblioteca C por baixo.

\paragraph{\emph{Bridge}: duas hierarquias que variam sozinhas.} É o padrão
cujo benefício é puramente aritmético, e o demo faz a conta:

\begin{lstlisting}[style=bt, language=bash,
    caption={Saída de \texttt{dist/bin/bridge}, seção 1: a explosão evitada.}]
Sem Bridge, com heranca nos dois eixos:
    3 formas x 3 renderizadores = 9 classes
    acrescentar 1 renderizador  = +3 classes
    acrescentar 1 forma         = +3 classes

Com Bridge:
    3 formas + 3 renderizadores = 6 classes
    acrescentar 1 renderizador  = +1 classe
    acrescentar 1 forma         = +1 classe

A conta piora rapido: 5x5 seria 25 contra 10.
\end{lstlisting}

O produto vira soma. E o demo mostra que os dois lados variam de verdade:
\texttt{redimensionar()} é do lado da \emph{abstração} e nenhum renderizador
soube dele; a troca de SVG para ASCII é do lado da \emph{implementação} e
nenhuma forma soube dela.

\begin{regra}[\emph{Bridge} \emph{vs.} \emph{Strategy}: a mesma forma, escalas diferentes]
Os dois têm diagrama idêntico --- um objeto guarda uma referência a uma
interface e delega. A diferença: no \emph{Strategy}, o outro lado é
\textbf{um método} (um algoritmo trocável); no \emph{Bridge}, o outro lado é
\textbf{uma hierarquia inteira} (uma família de operações que forma um
\emph{backend}). E há um caso degenerado famoso: o idioma \emph{pimpl}
é um \emph{Bridge} com exatamente uma implementação, usado
não para variar mas para esconder --- e para cortar tempo de compilação.
\end{regra}

\paragraph{\emph{Composite}: a árvore que o cliente trata como folha.} O demo
monta uma árvore de diretórios e chama a mesma função com um arquivo e com a
raiz inteira. A recursão está encapsulada em \texttt{Pasta::tamanho()}, e a
remoção de uma subárvore não tem um único \texttt{delete}: o
\classe{std::unique\_ptr} destrói tudo em cascata.

A decisão de design que o padrão força é onde colocar \texttt{adicionar()}, e o
subprojeto escolhe o lado menos ``uniforme'' de propósito:

\begin{table}[htbp]
\centering
\caption{As duas opções que o próprio GoF apresenta para o \emph{Composite}, e
         por que o subprojeto escolheu a segunda.}
\label{tab:composite-opcoes}
\footnotesize
\begin{tabular}{@{}lll@{}}
\toprule
\textbf{\texttt{adicionar()} fica\ldots} & \textbf{Uniformidade} & \textbf{Custo} \\
\midrule
na interface \texttt{No}, para todos & máxima & \texttt{Arquivo::adicionar()} viola LSP ou falha em silêncio \\
só no composto \texttt{Pasta}        & parcial & o cliente precisa do tipo \texttt{Pasta} para compor \\
\bottomrule
\end{tabular}
\end{table}

O argumento decisivo é prático: quase toda composição acontece na
\emph{construção} da árvore, onde os tipos concretos já são conhecidos. A
uniformidade que importa é a das operações de \emph{leitura} --- e essa a
segunda opção preserva inteira, sem obrigar \texttt{Arquivo} a prometer algo que
não sabe cumprir.

\begin{armadilha}[um bug real encontrado escrevendo o \emph{Composite}]
A primeira versão do demo declarava
\texttt{virtual void imprimir(int nivel = 0) const = 0;} --- e
\texttt{raiz->imprimir()} \textbf{não compilou}. O motivo: o \texttt{override}
em \texttt{Pasta} não repetiu o argumento padrão, e argumento padrão é escolhido
pelo tipo \textbf{estático}, não pelo dinâmico. A correção é o idioma NVI
(\cref{sec:oop}):
\begin{lstlisting}[style=bt, language=C++, frame=none,
                   xleftmargin=0pt, xrightmargin=0pt]
void imprimir() const { imprimir_em(0); }   // NAO-virtual, com o default
virtual void imprimir_em(int) const = 0;    // virtual, sem default
\end{lstlisting}
Aqui o NVI não é questão de estilo: é o que compila. Vale registrar que este bug
foi encontrado \emph{porque o alvo foi construído} --- um padrão bonito num
diretório que ninguém compila não ensina nada.
\end{armadilha}

\paragraph{\emph{Decorator}: composição em vez de herança.} O problema
--- acrescentar compressão, criptografia e \emph{logging} a um \emph{stream} em
\emph{qualquer combinação} --- resolvido por herança exigiria uma classe por
combinação. Com decoradores, cada um envolve o anterior:

\begin{lstlisting}[style=bt, language=C++,
    caption={Composição em tempo de execução: a ordem das camadas é a ordem
             das chamadas.}]
auto stream = std::make_unique<LoggingDecorator>(
    std::make_unique<CompressionDecorator>(
        std::make_unique<EncryptionDecorator>(
            std::make_unique<FileStream>("arquivo.bin")
        )
    )
);
stream->write("dados confidenciais"); // log -> comprime -> criptografa -> arquivo
\end{lstlisting}

\paragraph{\emph{Facade}: a ordem, escrita num lugar só.} O demo simula um
sistema de \emph{build} com seis subsistemas --- resolvedor de dependências,
compilador, cache, \emph{linker}, testes e empacotador. Sem a fachada, o cliente
escreve quinze linhas e precisa saber que \texttt{resolver} vem antes de
\texttt{compilar}, que a saída de \texttt{compilar} alimenta \texttt{linkar}, e
quais são os \emph{defaults} razoáveis de cada um. Esse conhecimento fica
duplicado em cada cliente. Com a fachada, é uma chamada.

O ponto que distingue a \emph{Facade} de uma camada de abstração é que
\textbf{o subsistema continua acessível}: quem precisa de um caso incomum
chama os subsistemas diretamente. A fachada otimiza o \emph{caso comum}, não
todos os casos --- e é por isso que ela não ``esconde'' nada.

\begin{armadilha}[a \emph{Facade} que cresce tem nome: \emph{God Object}]
Uma fachada que vira a \emph{única} porta de entrada, com quarenta métodos e
estado próprio, deixou de simplificar e passou a acoplar. O sinal de alarme é
objetivo: \textbf{ela acumular estado que os subsistemas não têm}. É a mesma
patologia da \emph{god class} do SRP (\cref{sec:solid}), e a mesma que ameaça o
\emph{Mediator}.
\end{armadilha}

\paragraph{\emph{Flyweight}: o único padrão do catálogo cujo ganho se mede.} A
motivação é puramente de memória, e a divisão que define o padrão é entre estado
\textbf{intrínseco} (compartilhável e imutável: o desenho da letra
\texttt{'a'}, a fonte, o tamanho, o \emph{bitmap} rasterizado) e
\textbf{extrínseco} (único por instância: a posição e a cor). Um documento de
cem mil caracteres precisa de cem mil \emph{posições}, mas de cerca de cem
glifos distintos.

\begin{lstlisting}[style=bt, language=bash,
    caption={Saída de \texttt{dist/bin/flyweight}, seção 3 --- a economia, medida
             no próprio binário deste repositório.}]
documento de 12200 caracteres:

COM flyweight:
  glifos distintos     :         21
  memoria dos glifos   :       5019 B
  memoria do documento :     292800 B  (12200 x 24 B)
  TOTAL                :     297819 B

SEM flyweight:
  objetos completos    :      12200
  TOTAL                :    3111000 B

economia: 10.4x menos memoria (3111000 B -> 297819 B)

sizeof(CaractereNaTela) =  24 B  (ponteiro + x,y + cor)
sizeof(CaractereGordo)  =  96 B  + as duas strings na heap
\end{lstlisting}

A assinatura do glifo é o que torna o compartilhamento possível:
\texttt{std::string desenhar(int x, int y, const std::string\& cor) const} ---
\texttt{x}, \texttt{y} e \texttt{cor} são \emph{parâmetros}, não membros. Um
objeto que guardasse a posição não poderia ser reusado por ninguém.

\begin{armadilha}[cinco condições sem as quais o \emph{Flyweight} não funciona]
O demo lista as cinco, e a segunda é a menos óbvia.
\textbf{(1)} O objeto compartilhado tem que ser \textbf{imutável} --- por isso a
fábrica devolve \texttt{const Glifo\&} e o glifo não tem \emph{setters}.
\textbf{(2)} O contêiner da fábrica não pode \textbf{invalidar referências}: o
demo usa \classe{std::deque} e não \classe{std::vector}, porque o
\texttt{vector} realoca ao crescer e transforma em pendurado todo ponteiro já
entregue (\cref{sec:stl}).
\textbf{(3)} A fábrica precisa sobreviver a todos os clientes.
\textbf{(4)} Em código concorrente, a \emph{inserção} precisa de sincronização;
a leitura dos glifos já criados é segura por serem imutáveis.
\textbf{(5)} Não compensa se o estado intrínseco for pequeno --- um ponteiro
custa 8 bytes, e compartilhar um objeto de 4 bytes \emph{gasta} memória. A
conta é: compensa se \texttt{sizeof(intrínseco)} $>$ \texttt{sizeof(ptr)} e
houver repetição real.
\end{armadilha}

O inverso interessante mora na biblioteca padrão: o \emph{SSO} da
\classe{std::string} --- medido em \projeto{strings} --- evita a indireção justamente para
objetos \emph{pequenos} --- pelo motivo exato da quinta condição.

\paragraph{\emph{Proxy}: mesma interface, acesso controlado.} O demo apresenta
três variedades --- \emph{proxy} virtual (adia a criação cara, usando
\texttt{mutable}), \emph{proxy} de proteção (verifica permissão antes de
repassar) e \emph{smart proxy} (registra, mede, faz \emph{cache}). E a pergunta
que o separa do \emph{Decorator} é única: \emph{o objeto real é sempre
chamado?} Se sim, é \emph{Decorator}; se ele pode \textbf{não} repassar ---
negar, adiar, servir do \emph{cache} --- é \emph{Proxy}. No demo, o
\emph{proxy} de proteção negou a chamada, e o \emph{proxy} de instrumentação
empilhado nem chegou a contá-la.

Vale notar que os \emph{smart pointers} da biblioteca padrão são \emph{proxies}
no sentido estrito, e o \classe{std::shared\_ptr} é o exemplo mais visitado do
catálogo sem que ninguém o chame assim.

% -----------------------------------------------------------------------------
\subsubsection{Os onze comportamentais}
% -----------------------------------------------------------------------------

Os comportamentais tratam de \emph{como os objetos se comunicam}. Quatro deles
já foram vistos de outro ângulo neste livro, e vale explicitar as ligações antes
dos novos.

\paragraph{\emph{Strategy} é OCP por injeção.} Cada algoritmo é um objeto
trocável em tempo de execução, e o contexto que o usa nunca é editado. É
literalmente o Demo~2 da \cref{sec:solid} com outro nome --- e a coincidência
não é acidental: o princípio descreve a propriedade, o padrão descreve a
realização.

\paragraph{\emph{Command} transforma ação em objeto.} O par
\texttt{execute()}/\texttt{undo()} é o que permite pilha de desfazer, fila de
tarefas e \emph{log} de operações. É a mesma reificação que a
\cref{sec:concurrency} usa para enfileirar trabalho numa \emph{thread pool} ---
com a diferença de que lá o objeto é um \classe{std::function}.

\paragraph{\emph{State} substitui o \texttt{switch} sobre estado.} Cada estado é
um objeto, e a transição é o próprio estado trocando o ponteiro do contexto ---
o exemplo do subprojeto é uma conexão TCP passando de \texttt{Listen} para
\texttt{Established}. O ganho sobre a máquina de estados com \texttt{enum} é o
mesmo do OCP: acrescentar um estado é acrescentar um arquivo.

\paragraph{\emph{Chain of Responsibility} e \emph{Observer}.} O primeiro
encadeia \emph{handlers} --- autenticação, limite de taxa, validação, negócio
--- e cada elo decide se processa ou repassa; o segundo notifica $N$ assinantes
a partir de um evento. São os dois padrões mais presentes em código de rede, e o
\cref{sec:network} os reencontra sem os nomear.

\paragraph{\emph{Template Method}: o que se protege é a \emph{ordem}.} O
esqueleto do algoritmo mora na base, num método \textbf{não virtual}, e os
derivados preenchem os buracos. A distinção de três níveis é o que faz o padrão
funcionar:

\begin{table}[htbp]
\centering
\caption{Os três tipos de passo do \emph{Template Method}, como o demo
         \texttt{dist/bin/template\_method} os classifica.}
\label{tab:template-method}
\footnotesize
\begin{tabular}{@{}lll@{}}
\toprule
\textbf{Tipo} & \textbf{Declaração} & \textbf{O derivado\ldots} \\
\midrule
obrigatório & \texttt{virtual \ldots{} = 0}        & \textbf{tem que} implementar \\
\emph{hook} & \texttt{virtual \ldots{} \{ padrão \}} & \textbf{pode} sobrescrever \\
fixo        & não virtual, \texttt{private}        & não pode mudar \\
\bottomrule
\end{tabular}
\end{table}

O demo quantifica o que cada variação precisou escrever: o
\texttt{PipelineCsv} implementou \textbf{2} métodos (só os obrigatórios), o
\texttt{PipelineJson} \textbf{3} (mais o \emph{hook} de validação) e o
\texttt{PipelineBanco} \textbf{6} (todos os \emph{hooks}). O passo
\texttt{ordenar()} é \texttt{private} e não virtual: ninguém o alcança. E
\texttt{executar()} é não virtual pelo mesmo motivo --- se um derivado pudesse
sobrescrevê-lo, a garantia de ordem, que é o padrão inteiro, evaporaria. É NVI
outra vez (\cref{sec:oop}).

\paragraph{\emph{Mediator}: a teia virando estrela.} O demo é um formulário com
seis componentes e regras cruzadas (``Enviar'' só ativa com e-mail válido
\emph{e} termos aceitos; marcar ``empresa'' revela o CNPJ e o torna
obrigatório). A conta que ele imprime é a razão de ser do padrão:

\begin{lstlisting}[style=bt, language=bash,
    caption={Saída de \texttt{dist/bin/mediator}, seção 1: quem conhece quem.}]
SEM mediador:
  email   -> botao, status
  termos  -> botao, status
  empresa -> cnpj, botao, status
  cnpj    -> botao, status
  6 componentes -> ate 30 referencias
  e acrescentar "aceito newsletter" toca varios deles

COM mediador:
  cada componente -> o mediador. So isso.
  6 referencias, e a regra num lugar so.
\end{lstlisting}

O ganho real não é a contagem: é que os componentes ficam \textbf{reutilizáveis}
--- um campo de texto que não conhece nenhum outro campo pode ir para qualquer
formulário.

\begin{armadilha}[o \emph{Mediator} carrega o defeito na própria definição]
Ao centralizar a interação, ele acumula toda a lógica que estava distribuída ---
e vira um \emph{God Object}. Os sinais: \texttt{notificar} virando um
\texttt{switch} de duzentas linhas; o mediador ganhando estado que não é de
ninguém; acrescentar um componente exigindo entender o mediador inteiro. As três
saídas que o demo registra: quebrar em mediadores menores por área; trocar o
\texttt{switch} por um \texttt{std::map<std::string, std::function<void()>>};
ou extrair as \emph{regras} para objetos próprios, deixando o mediador só
roteando.
\end{armadilha}

\paragraph{\emph{Memento}: salvar estado sem furar o encapsulamento.} São três
papéis --- o \textbf{originador} (que cria e restaura), o \textbf{memento}
(opaco) e o \textbf{zelador} (que guarda e nunca inspeciona). Em C++, a
opacidade não é convenção: é o sistema de tipos.

\begin{lstlisting}[style=bt, language=C++,
    caption={De \texttt{src/design-of-pattern/memento/}: a opacidade construída
             com classe aninhada privada e \texttt{friend}.}]
class Editor {
    class Estado {
        std::string conteudo_;   // privado
        std::size_t cursor_;     // privado
        Estado(...);             // construtor PRIVADO
        friend class Editor;     // so ele
    };
};
\end{lstlisting}

O que o zelador \emph{não} consegue fazer é o valor do padrão: ler
\texttt{conteudo\_}, escrever em \texttt{cursor\_} ou fabricar um
\texttt{Estado} do nada --- as três tentativas são erros de compilação. O que
ele pode, e é tudo de que precisa: guardar num \texttt{vector}, devolver ao
editor e ler os metadados públicos (rótulo, tamanho). O ganho concreto: o
originador pode \textbf{trocar a representação interna} sem que o histórico
sequer recompile.

\begin{regra}[\emph{Memento} \emph{vs.} \emph{Command}: a operação tem inversa barata?]
Se sim, use \emph{Command} e guarde a \emph{operação} --- custo $O(\Delta)$.
Inserir texto tem inversa óbvia (apagar). Se não --- um filtro de imagem
destrutivo ---, use \emph{Memento} e guarde o \emph{estado} inteiro, custo
$O(\text{estado})$. E os dois se combinam: um \emph{Command} que não consegue se
inverter guarda um \emph{Memento} do estado anterior, que é exatamente como
editores de imagem tratam filtros destrutivos. A distinção sutil é com o
\emph{Prototype}: o clone é um objeto \textbf{normal e usável}; o memento é
\textbf{opaco} e só serve para restaurar.
\end{regra}

\paragraph{\emph{Visitor}: despacho duplo, e o preço dele.} O \emph{Visitor}
resolve o problema de acrescentar operações a uma hierarquia fechada sem editar
as classes: cada nó tem \texttt{aceitar(Visitante\&)}, que chama de volta o
método específico do visitante --- daí ``duplo despacho''. Ele funciona, e o
demo o implementa por inteiro; mas o mesmo demo mede o que ele custa:

\begin{table}[htbp]
\centering
\caption{\emph{Visitor} clássico contra \texttt{variant} + \texttt{visit}, como
         o demo \texttt{dist/bin/visitor} compara na seção 6.}
\label{tab:visitor-variant}
\footnotesize
\begin{tabular}{@{}lll@{}}
\toprule
\textbf{Critério} & \textbf{\emph{Visitor}} & \textbf{\texttt{variant} + \texttt{visit}} \\
\midrule
operações novas               & fácil    & fácil \\
\textbf{tipos} novos          & difícil  & difícil \\
exaustividade verificada      & não      & \textbf{sim} \\
retorno com valor             & não      & \textbf{sim} \\
linhas de infraestrutura      & $\approx$ 40 & $\approx$ 3 \\
alocação por nó               & sempre   & opcional \\
tipos de bibliotecas alheias  & sim      & não \\
\bottomrule
\end{tabular}
\end{table}

Três defeitos do \emph{Visitor} clássico desaparecem com \texttt{visit}: a
operação passa a \textbf{devolver valor} (nada de acumular resultado num membro
do visitante), os nós deixam de precisar do método \texttt{aceitar}, e o
acoplamento circular entre nós e visitantes some. E há um ganho que o clássico
nunca teve: \texttt{visit} faz despacho \textbf{múltiplo}, sobre $N$
\texttt{variant} de uma vez --- estender o \emph{Visitor} para despacho triplo
seria uma explosão de métodos.

% -----------------------------------------------------------------------------
\subsubsection{As confusões que sobram}
% -----------------------------------------------------------------------------

Quase todo mal-entendido sobre padrões cai em um de oito pares. O
\texttt{README.md} do subprojeto os tabula, e cada demo dedica uma seção a
separar o seu padrão dos vizinhos.

\begin{table}[htbp]
\centering
\caption{Os oito pares que mais confundem, e a pergunta que separa cada um.}
\label{tab:confusoes}
\footnotesize
\begin{tabular}{@{}lp{0.56\textwidth}@{}}
\toprule
\textbf{Confusão} & \textbf{O que separa} \\
\midrule
\emph{Adapter} $\times$ \emph{Decorator} $\times$ \emph{Proxy} &
  \emph{Adapter} \textbf{muda} a interface; \emph{Decorator} \textbf{acrescenta} comportamento; \emph{Proxy} \textbf{controla} o acesso (pode não repassar) \\
\emph{Bridge} $\times$ \emph{Strategy} &
  o outro lado é uma \textbf{hierarquia} (\emph{Bridge}) ou \textbf{um método} (\emph{Strategy}) \\
\emph{Facade} $\times$ \emph{Mediator} &
  a \emph{Facade} \textbf{não acrescenta} lógica; o \emph{Mediator} \textbf{coordena} \\
\emph{Template Method} $\times$ \emph{Strategy} &
  o que se protege é a \textbf{ordem} dos passos ou o \textbf{algoritmo} inteiro? \\
\emph{Mediator} $\times$ \emph{Observer} &
  existe uma \textbf{regra} que depende de vários participantes? \\
\emph{Memento} $\times$ \emph{Command} &
  a operação tem \textbf{inversa barata}? \\
\emph{Memento} $\times$ \emph{Prototype} &
  o clone é \textbf{usável}; o memento é \textbf{opaco} \\
\emph{Visitor} $\times$ \texttt{variant} &
  a hierarquia é \textbf{aberta} (\emph{Visitor}) ou \textbf{fechada} (\texttt{variant})? \\
\bottomrule
\end{tabular}
\end{table}

% -----------------------------------------------------------------------------
\subsubsection{A versão moderna, padrão a padrão}
\label{sec:padroes-modernos}
% -----------------------------------------------------------------------------

Vários padrões do GoF descrevem soluções para limitações de linguagens dos anos
1990 --- linguagens sem \emph{closures}, sem uniões com tipo, sem inicialização
por nome. Em C++20, parte deles tem realização mais simples. Por isso os treze
padrões acrescentados na última fase do subprojeto trazem, além do padrão
clássico, uma seção \textbf{``versão moderna''} no próprio executável --- e é
dela que sai a \cref{tab:padroes-modernos}.

O ponto importante, e que a tabela sozinha não transmite: em \textbf{nenhum}
caso a versão moderna substitui a clássica sem trocar algo. A última coluna é a
que se deve ler primeiro.

\begin{table}[htbp]
\centering
\caption{A versão clássica, a moderna e o que se perde ao trocar. As perdas
         foram transcritas das seções ``versão moderna'' dos executáveis.}
\label{tab:padroes-modernos}
\scriptsize
\begin{tabular}{@{}lllp{0.30\textwidth}@{}}
\toprule
\textbf{Padrão} & \textbf{Clássica} & \textbf{Moderna} & \textbf{O que a moderna PERDE} \\
\midrule
\emph{Strategy} & hierarquia virtual & \classe{std::function} ou \emph{concept} & nome, estado e ciclo de vida do algoritmo \\
\emph{Command}  & classe com \texttt{execute()} & \emph{lambda} com \emph{init-capture} & o \texttt{undo()} pareado, que a \emph{lambda} não carrega \\
\emph{Visitor}  & duplo despacho com \texttt{accept()} & \classe{std::visit} sobre \texttt{variant} & tipos vindos de bibliotecas alheias \\
\emph{Iterator} & interface com \texttt{hasNext()} & iterador da STL, ou \emph{view} & --- (aqui a moderna ganha em tudo) \\
\emph{Singleton}& ponteiro estático + \emph{mutex} & \emph{Meyers singleton} & nada; mas o problema real continua (adiante) \\
\emph{Builder}  & \emph{builder} encadeado & \emph{designated initializers} & validação, campos obrigatórios, ordem livre, etapas \\
\emph{Prototype}& \texttt{clonar()} virtual & \texttt{variant} copiável, ou CRTP & a hierarquia deixa de ser aberta \\
\emph{Abstract Factory} & fábrica virtual por família & \texttt{struct} de \classe{std::function} & a garantia de consistência da família \\
\emph{Adapter}  & classe adaptadora & \emph{lambda} adaptadora & \gls{raii}, múltiplos métodos, estado próprio \\
\emph{Bridge}   & duas hierarquias & \texttt{struct} de \classe{std::function} & uma indireção por primitiva, em laço quente \\
\emph{Composite}& \texttt{No} virtual + filhos & \texttt{variant} recursivo & tipos de nó novos vindos de fora \\
\emph{Template Method} & \emph{hooks} virtuais & \texttt{struct} de \emph{callbacks} & a obrigatoriedade verificada em compilação \\
\emph{Flyweight}& fábrica + \texttt{const T\&} & \texttt{shared\_ptr<const T>} & a \textbf{garantia de unicidade} \\
\bottomrule
\end{tabular}
\end{table}

\paragraph{\emph{Builder} \emph{vs.} \emph{designated initializers}.} O C++20
permite \texttt{gerar(\{.titulo = "Vendas Q3", .formato = "json",
.casas\_decimais = 4\})}, e isso dá três das quatro coisas que se queria do
\emph{Builder}: cada campo tem nome no ponto de uso, os omitidos usam o
\emph{default} do \texttt{struct}, e o código extra é zero. O que \emph{não} dá
--- e é o que mantém o \emph{Builder} vivo --- são quatro coisas: não há
\textbf{validação} (o agregado aceita qualquer combinação), não há
\textbf{campos obrigatórios}, a \textbf{ordem não é livre} (em C++, os
designadores precisam seguir a ordem de declaração dos membros; em C, não), e
não há \textbf{construção em etapas} com lógica entre elas. A regra que sai
disso é simples: agregado sem invariantes, use \emph{designated init}; produto
com invariantes ou campos privados, use \emph{Builder}.

\paragraph{\emph{Template Method} \emph{vs.} \texttt{struct} de \emph{callbacks}.}
A versão moderna troca a hierarquia por um agregado de \classe{std::function},
e ganha o direito de montar os passos em tempo de execução, vindos de lugares
diferentes. Perde uma garantia que era gratuita: nada obriga a preencher o
\emph{callback} obrigatório. Um \classe{std::function} vazio \textbf{compila} e
lança \texttt{std::bad\_function\_call} em execução --- ao passo que o
\texttt{virtual \ldots{} = 0} tornava o esquecimento um erro de compilação.
Perde também o estado compartilhado entre os passos, que na versão clássica eram
membros protegidos da base e agora precisam virar capturas.

\paragraph{\emph{Bridge} e o custo do \emph{type erasure}.} Substituir a
hierarquia de renderizadores por um \texttt{struct} de \classe{std::function}
permite montar um \emph{backend} em execução, misturando primitivas de origens
diferentes. Mas cada \classe{std::function} é um \gls{typeerasure}, e cada
chamada custa uma indireção --- que num laço de renderização é caminho quente. A
medição de \texttt{src/lambdas/} registra \textbf{3,5$\times$} de diferença
entre a chamada direta e a passagem por \classe{std::function}
(\cref{sec:lambdas}). Para um \emph{Bridge} de renderização, a hierarquia
virtual costuma ganhar: a implementação é escolhida \emph{uma} vez na
inicialização, e depois são milhões de chamadas --- uma \gls{vtable} é mais
barata que $N$ \classe{std::function}. E se a escolha for conhecida em
compilação, um \emph{template} elimina os dois custos.

\paragraph{\emph{Flyweight} e o \classe{std::shared\_ptr}.} Trocar a fábrica por
\texttt{shared\_ptr<const Glifo>} ganha contagem de referência automática --- os
glifos sem uso podem ser liberados. Perde a única coisa que o padrão promete:
\textbf{nada impede criar dois \texttt{shared\_ptr} para glifos idênticos}. É a
fábrica que garante que ``mesmo estado intrínseco'' implica ``mesmo objeto''; o
ponteiro compartilhado, sozinho, não garante nada.

\paragraph{\emph{Singleton}: a versão moderna resolveu o problema errado.} O
\emph{Meyers singleton} é \emph{thread-safe} por garantia do padrão desde o
C++11 --- a inicialização de um \texttt{static} local é sincronizada pelo
compilador:

\begin{lstlisting}[style=bt, language=C++,
    caption={O \emph{Meyers singleton}: \emph{thread-safe} por garantia do padrão.}]
class Logger {
public:
    static Logger& instance() {
        static Logger inst; // thread-safe por garantia do C++11 (6.7)
        return inst;
    }

    Logger(const Logger&)            = delete;
    Logger& operator=(const Logger&) = delete;

private:
    Logger() = default;
};
\end{lstlisting}

\begin{armadilha}[o problema do \emph{Singleton} nunca foi a inicialização]
Resolver a inicialização preguiçosa e a segurança entre \emph{threads} não
resolve o que realmente pesa: estado global que persiste entre testes,
acoplamento oculto (qualquer código pode chamar \texttt{Logger::instance()}) e
violação direta do DIP --- a dependência não aparece na assinatura. É a mesma
crítica que se faz ao \emph{service locator} (\cref{sec:solid}). O uso legítimo
existe --- \emph{logger} de processo, configuração, registro de \emph{plugins}
(\cref{sec:plugin-registry}) --- mas a pergunta certa continua sendo se a
injeção por construtor não resolveria melhor.
\end{armadilha}

E o \emph{Null Object} é o complemento que torna o \emph{Singleton} dispensável
em teste: uma implementação que não faz nada, injetada no lugar da real. O ganho
é eliminar as verificações de \texttt{nullptr} do cliente --- o membro passa a
ser uma \emph{referência}, e referência nunca é nula. O demo mostra os dois
lados da moeda: com um \texttt{NullLogger}, um erro de validação é
\emph{silenciado} --- e ninguém percebe, porque não há verificação nenhuma.

\begin{regra}[\texttt{optional} para dados; \emph{Null Object} para comportamento]
\classe{std::optional} obriga o chamador a verificar --- e é isso que se quer
para um \emph{valor} que pode não existir (\cref{sec:vocabulary}). O \emph{Null
Object} faz o contrário: garante que a chamada seja segura para que ninguém
precise verificar --- e é isso que se quer para um \emph{serviço} opcional. Usar
um no lugar do outro produz ou verificações desnecessárias ou ausência
silenciosa de comportamento.
\end{regra}

\subsubsection{\emph{Interpreter}: onde ele ganha e onde perde}

O \emph{Interpreter} do subprojeto implementa a mesma gramática de expressões
aritméticas que o \projeto{parser} do \cref{sec:parser} resolve com Flex e
Bison, o que permite a comparação direta (\cref{tab:interpreter-bison}). A
precedência dos operadores sai naturalmente da estrutura das chamadas:
\texttt{expr} chama \texttt{term}, que chama \texttt{factor}.

\begin{table}[htbp]
\centering
\caption{\emph{Interpreter} escrito à mão contra gerador de \emph{parser}.}
\label{tab:interpreter-bison}
\footnotesize
\begin{tabular}{@{}lll@{}}
\toprule
 & \textbf{\emph{Interpreter} GoF} & \textbf{Flex/Bison} \\
\midrule
dependências           & zero --- C++ puro      & \texttt{flex} + \texttt{bison} instalados \\
AST                    & objeto reutilizável    & precisa ser implementada à parte \\
testabilidade          & cada nó isoladamente   & o \emph{parser} inteiro \\
gramáticas grandes     & trabalhoso             & ideal \\
recuperação de erros   & manual                 & robusta \\
quando usar            & DSL embutida, $<$ 20 regras & linguagem completa \\
\bottomrule
\end{tabular}
\end{table}

Acrescentar um operador novo ao \emph{Interpreter} são três passos --- classe
nova, função auxiliar, um caso no \emph{parser} --- e nenhum arquivo existente
quebra. É OCP em ação, e é também exatamente o eixo em que uma AST baseada em
\texttt{std::variant} \emph{perderia}: acrescentar um \emph{tipo} de nó a um
\texttt{variant} recompila todos os \texttt{visit} (\cref{sec:vocabulary}). É o
\emph{expression problem} da \cref{sec:expression-problem} aparecendo, pela
última vez, com nome e sobrenome.

% =============================================================================
\subsection{Entity Component System}
\label{sec:ecs}
% =============================================================================

O subprojeto \projeto{ecs} implementa do zero o padrão arquitetural dominante
nas \emph{engines} de jogo --- Unity, Unreal, Bevy. Ele entra na trilha logo
depois dos padrões porque é o \textbf{contraponto direto} ao \cref{sec:oop}:
composição em tempo de execução no lugar de herança.

\subsubsection{O problema}

Uma hierarquia de herança obriga a decidir, em tempo de compilação, quais
combinações de capacidades existem --- e o número de combinações cresce
exponencialmente:

\begin{lstlisting}[style=bt, language=C++,
    caption={A explosão combinatória, e a alternativa.}]
class EntidadeMovel   : public Entidade { Posicao pos; Velocidade vel; };
class EntidadeComSaude: public Entidade { Saude hp; };
class Jogador : public EntidadeMovel, public EntidadeComSaude { };  // e as outras?

Entity jogador = world.criar();
world.adicionar<Posicao>(jogador, 0.f, 0.f);
world.adicionar<Velocidade>(jogador, 5.f, 0.f);
world.adicionar<Saude>(jogador, 100, 100);
world.adicionar<Voando>(jogador);       // vira voadora em runtime, sem classe nova
\end{lstlisting}

\subsubsection{Os três pilares}

\begin{table}[htbp]
\centering
\caption{A separação radical entre identidade, dados e lógica.}
\label{tab:ecs-pilares}
\footnotesize
\begin{tabular}{@{}lll@{}}
\toprule
\textbf{Pilar} & \textbf{Responsabilidade} & \textbf{Neste subprojeto} \\
\midrule
\textbf{Entity}    & identificador único (só um número) & \texttt{uint32\_t} com reciclagem \\
\textbf{Component} & dados puros, sem lógica            & \texttt{Posicao}, \texttt{Velocidade}, \texttt{Saude} \\
\textbf{System}    & lógica que opera sobre componentes & \texttt{movimento()}, \texttt{colisao()}, \texttt{dano()} \\
\bottomrule
\end{tabular}
\end{table}

A entidade não é um objeto: é um número. Não tem métodos, não tem dados, não tem
tipo. Os componentes são \emph{aggregates} sem um único método --- inclusive os
de marcação, que não têm nem dados (\texttt{struct Jogador \{\};}) e existem só
para que a presença deles seja consultável. E os sistemas são funções livres.

\begin{porque}[por que a reciclagem de IDs importa]
Um jogo de longa duração cria e destrói projéteis a cada quadro. Sem
reciclagem, o espaço de IDs cresce indefinidamente até estourar o
\texttt{uint32\_t}; com reciclagem, o \texttt{EntityManager} devolve o ID
destruído a uma pilha e o reutiliza na próxima criação. O custo é que um ID
reutilizado pode ser confundido com o antigo --- razão pela qual ECS de
produção usam \emph{generational indices}, com um contador de geração no bit
alto do identificador.
\end{porque}

\subsubsection{\emph{Type erasure} para guardar tipos heterogêneos}

O problema técnico central é armazenar \texttt{ComponentStore<Posicao>},
\texttt{ComponentStore<Saude>} e todas as outras no mesmo mapa --- tipos
diferentes, sem base comum natural. A solução é o mesmo \gls{typeerasure} de
\texttt{std::function} e \texttt{std::any} (\cref{sec:lambdas,sec:vocabulary}):
uma interface não \emph{template}, e um \classe{std::type\_index} como chave.

\begin{lstlisting}[style=bt, language=C++,
    caption={A chave é o tipo, e o \texttt{static\_cast} de volta é seguro
             porque a chave garante qual é.}]
struct IComponentStore {                    // interface base, sem templates
    virtual ~IComponentStore() = default;
    virtual void remover(Entity e) = 0;
};

template<typename T>
class ComponentStore : public IComponentStore { /* ... */ };

std::unordered_map<std::type_index, std::unique_ptr<IComponentStore>> stores_;

template<typename T>
ComponentStore<T>& store_de() {
    auto key = std::type_index(typeid(T));
    return static_cast<ComponentStore<T>&>(*stores_[key]);
}
\end{lstlisting}

\subsubsection{A \emph{query}: onde as peças se encaixam}

A operação que define o ECS é a consulta ``todas as entidades que têm
\emph{estes} componentes''. E ela é uma aplicação direta de uma
\emph{fold expression} (\cref{sec:templates}):

\begin{lstlisting}[style=bt, language=C++,
    caption={A \emph{query} genérica, em \texttt{src/ecs/world.hpp}.}]
    template <typename... Ts, typename Func>
    void query(Func&& func) {
        // Itera sobre todas as entidades do primeiro tipo (menor iteração possível)
        // Depois verifica se têm os demais tipos
        auto& primeiro_store = store_de<first_t<Ts...>>();

        const auto& ents = primeiro_store.entidades();
        for (std::size_t i = 0; i < ents.size(); ++i) {
            Entity e = ents[i];
            if (tem_todos<Ts...>(e)) {
                func(e, get<Ts>(e)...);
            }
        }
    }
\end{lstlisting}

Com ela, um sistema inteiro cabe em quatro linhas, e é uma função pura testável
isoladamente:

\begin{lstlisting}[style=bt, language=C++,
    caption={Um sistema completo.}]
void movimento(World& w, float dt) {
    w.query<Posicao, Velocidade>([&](Entity, Posicao& p, Velocidade& v) {
        p.x += v.dx * dt;
        p.y += v.dy * dt;
    });
}
\end{lstlisting}

\subsubsection{O que se ganha e o que se perde}

\begin{table}[htbp]
\centering
\caption{ECS contra hierarquia de herança.}
\label{tab:ecs-heranca}
\footnotesize
\begin{tabular}{@{}lll@{}}
\toprule
\textbf{Critério} & \textbf{Herança} & \textbf{ECS} \\
\midrule
combinação de capacidades   & explosão combinatória & composição em execução \\
localidade de cache         & objetos dispersos     & componentes do mesmo tipo juntos \\
acrescentar comportamento   & classe nova           & \texttt{adicionar<T>()} em execução \\
paralelismo                 & difícil               & sistemas independentes \\
testabilidade               & precisa do objeto     & função pura sobre componentes \\
legibilidade do fluxo       & direta                & indireta (quem roda quando?) \\
\bottomrule
\end{tabular}
\end{table}

\begin{armadilha}[o ECS deste subprojeto ainda não é \emph{cache-friendly}]
A promessa de localidade de cache do ECS depende de os componentes ficarem
\emph{contíguos} em memória --- o que exige armazenamento em \texttt{vector}
denso com um mapa esparso de índices. O \texttt{ComponentStore} aqui usa
\texttt{std::unordered\_map<Entity, T>}, que espalha os componentes pela
\emph{heap} exatamente como a herança fazia. A arquitetura está certa e o
armazenamento é didático: trocá-lo por um \emph{sparse set} não muda uma linha
dos sistemas --- que é, aliás, a melhor demonstração de que a separação vale a
pena.
\end{armadilha}

% =============================================================================
\subsection{Plugins: DIP até o limite do sistema operacional}
\label{sec:plugin-registry}
% =============================================================================

O subprojeto \projeto{plugin\_registry} carrega bibliotecas compartilhadas em
tempo de execução com \texttt{dlopen}/\texttt{dlsym}, descobre os
\emph{plugins} disponíveis varrendo um diretório, e troca o algoritmo de
controle sem recompilar nada. É o \cref{sec:solid} levado ao extremo: a
dependência não só está invertida como \textbf{nem existe no binário}.

O contrato é uma interface virtual pura --- \texttt{IController}, com
\texttt{name()}, \texttt{reset()} e \texttt{compute()} --- e o que atravessa a
fronteira são apenas \emph{aggregates} simples (\texttt{State},
\texttt{Action}).

\subsubsection{O que é, exatamente, carregar código em tempo de execução}

Antes do código, os termos. Um \textbf{objeto compartilhado} (\emph{shared
object}, arquivo \texttt{.so} no Linux, \texttt{.dll} no Windows,
\texttt{.dylib} no macOS) é um binário que contém código de máquina e uma
\textbf{tabela de símbolos}: uma lista que diz, para cada função exportada, em
que deslocamento dentro do arquivo ela começa. Nada nele é endereço absoluto ---
é código \emph{relocável}, compilado com \texttt{-fPIC}, para poder ser
carregado em qualquer posição da memória do processo.

Há duas maneiras de usar um \texttt{.so}, e a diferença entre elas é o assunto
deste subprojeto (\cref{tab:ligacao}). Na \textbf{ligação dinâmica comum}, o
\emph{linker} anota no executável ``preciso de \texttt{libfmt.so}'', e quem
carrega é o \emph{loader} do sistema, antes de o \texttt{main} começar. Na
\textbf{ligação dinâmica explícita}, ninguém anota nada: o programa chama
\texttt{dlopen} em tempo de execução, com um caminho que pode ter vindo de um
arquivo de configuração, e o arquivo pode não existir na hora da compilação.

\begin{table}[htbp]
\centering
\caption{As duas formas de usar uma biblioteca compartilhada. Só a segunda
         permite que o \emph{plugin} seja escrito depois do \emph{host}.}
\label{tab:ligacao}
\footnotesize
\begin{tabular}{@{}lll@{}}
\toprule
 & \textbf{Ligação dinâmica comum} & \textbf{\texttt{dlopen} explícito} \\
\midrule
quem carrega            & o \emph{loader}, antes do \texttt{main} & o seu código, quando quiser \\
o nome do \texttt{.so}  & fixo no executável            & uma \texttt{std::string} qualquer \\
símbolo faltando        & o programa nem inicia          & exceção tratável \\
o arquivo precisa existir & na compilação                & só na hora do \texttt{dlopen} \\
resolução de símbolo    & pelo tipo, no \emph{link}      & \textbf{por nome literal}, em execução \\
\bottomrule
\end{tabular}
\end{table}

A última linha é a que dita todo o desenho que vem a seguir: \texttt{dlsym}
procura uma \emph{string}. Não há verificação de tipos --- o que atravessa a
fronteira é um nome e um \texttt{void*}.

\begin{porque}[por que só \emph{aggregates} simples cruzam a fronteira]
\texttt{State} e \texttt{Action} são \texttt{struct} de \texttt{double}, sem
\gls{vtable}, sem alocação escondida, sem \classe{std::string} dentro. Isso é
deliberado, e o cabeçalho do subprojeto documenta o porquê: \emph{host} e
\emph{plugin} são duas unidades compiladas separadamente, possivelmente por
compiladores ou versões diferentes de biblioteca padrão. Um
\classe{std::string} passado por valor entre elas depende de os dois lados
concordarem sobre o \emph{layout} interno da \texttt{string} --- e essa
concordância é a \gls{abi}, que nem o padrão da linguagem garante. Um
\texttt{struct} de \texttt{double} tem \emph{layout} ditado pela ABI da
plataforma (System V AMD64, no Linux), e essa, sim, é estável há décadas.
\end{porque}

\subsubsection{Os símbolos exportados}

\begin{lstlisting}[style=bt, language=C++,
    caption={Os três símbolos que todo \emph{plugin} precisa exportar.}]
extern "C" const char* plugin_name() { return "pid"; }

extern "C" plugin::IController* create_controller() {
    return new PIDController();
}

extern "C" void destroy_controller(plugin::IController* p) {
    delete p;
}
\end{lstlisting}

O \texttt{extern "C"} é obrigatório: \texttt{dlsym} busca por nome literal, e a
decoração de nomes do C++~\cite{itaniumabi} transformaria
\texttt{create\_controller} em algo que ninguém digitaria. Note que o
\emph{plugin} exporta funções livres, não a classe --- a classe permanece
invisível, e o \emph{host} só conhece a interface.

Como os três símbolos são sempre os mesmos, o subprojeto os empacota numa macro,
e o cabeçalho \texttt{icontroller.hpp} explica, no próprio comentário, por que
uma macro e não algo mais civilizado:

\begin{lstlisting}[style=bt, language=C++,
    caption={De \texttt{src/plugin\_registry/icontroller.hpp}: um \emph{plugin}
             inteiro se registra com uma linha.}]
#define REGISTER_CONTROLLER(ClassName, name_str)                        \
    extern "C" __attribute__((visibility("default")))                   \
    const char* plugin_name() { return name_str; }                      \
    extern "C" __attribute__((visibility("default")))                   \
    plugin::IController* create_controller() { return new ClassName(); }\
    extern "C" __attribute__((visibility("default")))                   \
    void destroy_controller(plugin::IController* p) { delete p; }

// no plugin, uma linha:
REGISTER_CONTROLLER(PIDController, "pid")
\end{lstlisting}

As três alternativas foram descartadas por motivos concretos, e vale enunciá-los
porque é raro ver uma macro justificada de forma defensável. Um \emph{template}
não serve: \emph{templates} têm decoração de nomes, e \texttt{dlsym} não
acharia nada. Uma função normal de auto-registro não serve: exigiria que o
\emph{plugin} se registrasse em algum ponto de inicialização estática --- com
todas as armadilhas de ordem de inicialização já vistas na \cref{sec:solid}. A
macro resolve no pré-processador: zero custo em execução, zero infraestrutura,
zero risco de esquecer um dos três símbolos.

O \texttt{\_\_attribute\_\_((visibility("default")))} é o complemento
necessário: com \texttt{-fvisibility=hidden} --- que é o padrão recomendado
para bibliotecas compartilhadas, porque encolhe a tabela de símbolos e acelera o
carregamento --- os símbolos ficam \emph{invisíveis} por omissão, e o
\texttt{dlsym} não acharia justamente aqueles três. O atributo abre exceção
para eles, e apenas para eles.

\subsubsection{\texttt{dlopen} e a ordem certa de checar erro}

\begin{lstlisting}[style=bt, language=C++,
    caption={As flags, e o motivo de \texttt{dlerror()} vir antes.}]
void* handle = dlopen(path.c_str(), RTLD_NOW | RTLD_LOCAL);
if (!handle) throw std::runtime_error("dlopen falhou: " + std::string(dlerror()));

dlerror();                                     // limpa erro anterior
auto fn = reinterpret_cast<FnPluginName>(dlsym(handle, "plugin_name"));
if (const char* err = dlerror()) {             // e só ENTÃO checa
    dlclose(handle);
    throw std::runtime_error("plugin_name não encontrado: " + std::string(err));
}
\end{lstlisting}

\texttt{RTLD\_NOW} resolve todos os símbolos no carregamento, fazendo a falha
acontecer cedo em vez de na primeira chamada; \texttt{RTLD\_LOCAL} impede que os
símbolos do \emph{plugin} vazem para os outros. As alternativas existem e têm
uso legítimo (\cref{tab:rtld}) --- mas as duas escolhas do subprojeto são as
que se deve tomar por omissão.

\begin{table}[htbp]
\centering
\caption{As quatro \emph{flags} de \texttt{dlopen} que importam, e quando a
         alternativa se justifica.}
\label{tab:rtld}
\footnotesize
\begin{tabular}{@{}lp{0.34\textwidth}p{0.30\textwidth}@{}}
\toprule
\textbf{\emph{Flag}} & \textbf{O que faz} & \textbf{Quando escolher} \\
\midrule
\texttt{RTLD\_NOW}   & resolve todos os símbolos ao carregar & padrão: falha cedo, com mensagem \\
\texttt{RTLD\_LAZY}  & resolve cada símbolo na primeira chamada & \texttt{.so} enorme do qual se usa pouco \\
\texttt{RTLD\_LOCAL} & os símbolos não vazam para outros \texttt{.so} & padrão: isola os \emph{plugins} \\
\texttt{RTLD\_GLOBAL}& os símbolos ficam globais & quando um \emph{plugin} depende de outro \\
\bottomrule
\end{tabular}
\end{table}

\begin{armadilha}[\texttt{RTLD\_LAZY} adia o erro para o pior momento possível]
Com \texttt{RTLD\_LAZY}, um \emph{plugin} a que falte um símbolo carrega sem
reclamar e mata o processo horas depois, na primeira vez que aquele caminho de
código for percorrido --- tipicamente em produção, e nunca no teste. É a mesma
categoria de erro dos alvos que não se constroem, discutida no fim desta seção:
o problema não é a falha, é o \emph{atraso} entre a causa e o sintoma.
\end{armadilha}

\begin{armadilha}[\texttt{dlsym} pode devolver \texttt{nullptr} legitimamente]
Testar o retorno de \texttt{dlsym} contra \texttt{nullptr} é errado: um símbolo
pode existir e \emph{valer} \texttt{nullptr}. A forma correta é chamar
\texttt{dlerror()} antes, para limpar qualquer erro pendente, e conferir
\texttt{dlerror()} depois. Como o repositório documenta no próprio código, é
esta a ordem, e não a intuitiva.
\end{armadilha}

\subsubsection{O \emph{deleter} customizado, e por que ele é crítico}

Este é o ponto em que o subprojeto deixa de ser sobre \emph{plugins} e passa a
ser sobre \gls{abi}. O objeto foi alocado \emph{dentro} do \texttt{.so}, com o
alocador do \texttt{.so}. Se o \emph{host} fizer \texttt{delete}, ele usará o
alocador \emph{dele} --- e se os dois forem diferentes, o resultado é corrupção
de \emph{heap}.

\begin{lstlisting}[style=bt, language=C++,
    caption={O \emph{deleter} como parâmetro de \emph{template}, em
             \texttt{src/plugin\_registry/plugin\_registry.hpp}.}]
    std::unique_ptr<IController, FnDestroyController>
    create(const std::string& name) {
        // ...
        // O segundo argumento do unique_ptr é o deleter:
        // em vez de `delete ptr`, ele chama `ph.destroy(ptr)`.
        return { raw, ph.destroy };
    }
\end{lstlisting}

\begin{regra}[quem aloca, destrói]
Um objeto criado dentro de uma biblioteca dinâmica deve ser destruído por ela.
O \texttt{unique\_ptr} com \emph{deleter} customizado
(\cref{sec:smart-pointers}) é o instrumento exato para isso: a regra deixa de
depender da disciplina de quem escreve o código de chamada e passa a estar no
\emph{tipo} do ponteiro. É a mesma técnica do \texttt{unique\_ptr<FILE,
decltype(closer)>}, aplicada a um problema muito mais caro de diagnosticar.
\end{regra}

\subsubsection{O registro, e a ordem em que tudo morre}

O \texttt{ControllerRegistry} é um \emph{Meyers singleton}: uma função com
\texttt{static} local, exatamente como a \cref{sec:padroes-modernos} descreve.
Aqui o uso é dos legítimos --- há \emph{handles} do sistema operacional abertos,
e faz sentido que exista um dono só ---, e a classe declara isso no tipo:
construtor de cópia e atribuição são \texttt{= delete}.

O destrutor fecha todos os \emph{handles} com \texttt{dlclose}, e é aí que mora
a armadilha mais cara deste subprojeto.

\begin{armadilha}[\texttt{dlclose} antes do último \texttt{unique\_ptr} é código morto sendo executado]
Quando \texttt{dlclose} descarrega um \texttt{.so}, o sistema operacional
\textbf{desmapeia as páginas} em que aquele código estava. Todos os ponteiros
para dentro dele --- ponteiros de função, ponteiros para as \gls{vtable} das
classes do \emph{plugin}, o próprio \emph{deleter} guardado dentro do
\texttt{unique\_ptr} --- passam a apontar para memória não mapeada. Destruir um
\texttt{unique\_ptr<IController, FnDestroyController>} depois disso não é ``usar
um ponteiro inválido'': é saltar para um endereço que não existe mais, e o
processo morre com \texttt{SIGSEGV} num \emph{stack trace} que não menciona
nenhuma linha do seu código.

A regra que evita isso é de tempo de vida, não de sintaxe: \textbf{todos os
\texttt{unique\_ptr} devem ser destruídos antes do registro}. O comentário no
destrutor de \texttt{plugin\_registry.hpp} registra exatamente essa
dependência --- e ela é a razão de o \texttt{main} criar os controladores dentro
de funções de escopo curto (\texttt{run\_episode}), em vez de guardá-los em
variáveis globais.
\end{armadilha}

\subsubsection{Onde estão os \emph{plugins}? A resolução de caminho}

Um detalhe que parece trivial e não é. O \emph{host} precisa saber em que
diretório procurar os \texttt{.so} --- e a resposta ingênua, um caminho relativo
como \texttt{"plugins"}, resolve contra o \textbf{diretório de trabalho} do
processo, não contra o local do executável. O mesmo binário passa a funcionar ou
não dependendo de onde foi invocado: roda de dentro do \texttt{dist/}, falha se
chamado por caminho absoluto de qualquer outro lugar.

A saída portável no Linux é \texttt{/proc/self/exe}: um \emph{symlink} que o
núcleo mantém apontando para o binário do próprio processo. Lê-lo devolve o
caminho real, já com todos os \emph{symlinks} resolvidos.

\begin{lstlisting}[style=bt, language=C++,
    caption={De \texttt{src/plugin\_registry/main.cpp}: descobrir o próprio
             diretório, sem depender do \texttt{CWD}.}]
static std::filesystem::path diretorio_do_executavel() {
    std::error_code ec;
    auto exe = std::filesystem::read_symlink("/proc/self/exe", ec);
    return ec ? std::filesystem::path{} : exe.parent_path();
}
\end{lstlisting}

Com isso, a busca vira uma lista de candidatos com prioridade explícita, e cada
um deles corresponde a um \emph{layout} real do repositório:

\begin{lstlisting}[style=bt, language=C++,
    caption={A ordem dos candidatos --- e o argumento de linha de comando
             vencendo todos.}]
if (argc > 1) return argv[1];                  // 1. explicito ganha sempre
candidatos.push_back(exe_dir / ".." / "plugins"); // 2. instalado: dist/bin -> dist/plugins
candidatos.push_back(exe_dir / "plugins");        // 3. build: ao lado do exe
candidatos.emplace_back("plugins");               // 4. ultimo recurso: o CWD
\end{lstlisting}

E quando nenhum serve, o programa não estoura uma exceção não tratada: imprime
os caminhos que \emph{tentou}, ensina a sintaxe do argumento explícito e devolve
\texttt{1}. É uma diferença pequena de código e enorme de diagnóstico --- a
mensagem contém a informação que o usuário precisa para consertar sozinho.

\begin{regra}[três formas de achar recursos, e a que se deve preferir]
\texttt{argv[0]} não serve: ele contém o que o \emph{shell} digitou, que pode
ser um nome relativo ou nem existir. Variáveis de ambiente
(\texttt{MEUAPP\_PLUGIN\_DIR}) servem como \emph{override} e são o que se
oferece ao usuário avançado. Mas o \emph{default} deve ser derivado do
executável --- \texttt{/proc/self/exe} no Linux,
\texttt{GetModuleFileName} no Windows, \texttt{\_NSGetExecutablePath} no
macOS --- porque é a única fonte que não depende de como o programa foi
invocado.
\end{regra}

\subsubsection{O ciclo completo}

A descoberta automática fecha o ciclo: o registro varre o diretório, tenta
carregar cada \texttt{.so}, e --- detalhe que importa em produção --- captura a
exceção por \emph{plugin}, de modo que um \emph{plugin} defeituoso não derruba
os demais. O resultado é a troca de algoritmo em tempo de execução, sem
recompilação:

\begin{lstlisting}[style=bt, language=bash,
    caption={Saída real de \texttt{dist/bin/plugin\_registry}, executado sem
             argumento nenhum: o caminho foi deduzido de
             \texttt{/proc/self/exe}.}]
=== Plugin Registry Demo ===

Varrendo '/home/lima/Projects/cpp_academy/dist/plugins'...
[registry] plugin 'pid' carregado de .../dist/plugins/pid.so
[registry] plugin 'lqr' carregado de .../dist/plugins/lqr.so
2 plugin(s) carregado(s).

=== Troca de controlador em runtime ===
(mesma interface, implementacao diferente, zero recompilacao)

  [pid]   alt= 2900.0 ft  spd= 88.0 kts  pitch=  0.0  | elev=+1.000  thr=0.750
  [lqr]   alt= 2900.3 ft  spd= 88.0 kts  pitch=  2.0  | elev=-0.011  thr=0.710
\end{lstlisting}

As duas últimas linhas são o subprojeto inteiro em duas medições: o
\textbf{mesmo estado} entra nos dois controladores e sai uma ação diferente de
cada um. O PID satura o profundor em \texttt{+1.000}; o LQR, com a mesma
entrada, pede \texttt{-0.011}. O \emph{host} não sabe qual é qual, nunca soube
o nome das classes, e não foi recompilado entre as duas chamadas.

\subsubsection{A lição que só apareceu quando os alvos foram construídos}

Vale fechar este capítulo com um resultado de método, e não de projeto. Até uma
revisão recente, boa parte dos alvos deste repositório estava com o
\texttt{subdir()} comentado no \texttt{src/meson.build} --- inclusive os
``pesados'' (Qt, SFML, Bullet, ONNX, JSBSim, AWS, RabbitMQ, SOCI). Hoje os
\textbf{57} subprojetos são construídos, produzindo \textbf{98} binários
instalados em \texttt{dist/bin}, com zero avisos sob
\texttt{-Wall -Wextra -Wpedantic}.

Habilitar os alvos não foi uma operação administrativa: ela \emph{encontrou
bugs}, e todos estavam exatamente nos alvos que ninguém compilava.

\begin{table}[htbp]
\centering
\caption{Bugs reais encontrados ao habilitar alvos que estavam desativados. A
         coluna da direita é o que cada um ensina --- e nenhum deles teria
         aparecido em revisão de código.}
\label{tab:bugs-latentes}
\footnotesize
\begin{tabular}{@{}lp{0.52\textwidth}@{}}
\toprule
\textbf{Onde} & \textbf{O que estava errado} \\
\midrule
\texttt{asm\_course/07} & um \texttt{subq \$8, \%rsp} quebrava o alinhamento de 16 bytes que a ABI System V exige no \texttt{call}; o \emph{crash} aparecia \emph{dentro} do \texttt{fmt}, quatro quadros acima \\
\texttt{asm\_course/10} & \texttt{RDI}/\texttt{RSI}/\texttt{RCX} declarados como entrada pura num \emph{asm} \emph{inline}, mas \texttt{rep movsb} os modifica --- o GCC acreditou na promessa e reusou um registrador corrompido \\
\texttt{http-client}    & \texttt{std::string(n, '─')}: U+2500 tem \textbf{três} bytes, e entre aspas simples vira \emph{multi-character constant} truncada; e \texttt{40 - title.size()} é aritmética sem sinal, que dá a volta \\
\texttt{aws/s3}         & faltava \texttt{\#include <fmt/core.h>} --- latente porque o alvo nunca era compilado \\
\texttt{rabbitmq}       & a API do AMQP-CPP mudou: os \emph{callbacks} de ciclo de vida são métodos virtuais do \emph{handler}, não do \texttt{TcpConnection} \\
\texttt{database}       & o demo não era idempotente: a segunda execução batia numa restrição \texttt{UNIQUE} \\
\bottomrule
\end{tabular}
\end{table}

\begin{regra}[um alvo que não é construído é um alvo que apodrece]
Quatro \texttt{\#include} faltando e dois bugs de \gls{abi} estavam
\emph{exatamente} nos alvos desabilitados --- e não é coincidência. Um alvo
comentado não é ``código pausado'': é código que continua envelhecendo enquanto
as bibliotecas mudam de API, os cabeçalhos param de ser incluídos
transitivamente e o compilador aprende avisos novos, sem que nada nunca o diga.
O corolário é organizacional: o mesmo raciocínio que levou o diretório
\texttt{tests/} da raiz a ser \textbf{removido} --- ele estava 100\%
comentado e nunca construiu nada, dando a \emph{impressão} de cobertura sem
oferecer nenhuma. Os testes hoje moram junto do código que testam
(\texttt{src/testing/} com 5 alvos e 41 casos, \texttt{src/meson/} com 1 alvo,
\texttt{src/cmake/} com CTest e 11 casos), e \texttt{meson test} roda 6 alvos,
todos passando. Um teste desabilitado é pior que teste nenhum, porque mente.
\end{regra}

E o alinhamento entre esta lição e o assunto do capítulo é direto: as fronteiras
de que SOLID, os padrões e o registro de \emph{plugins} tratam só cumprem o que
prometem se forem \emph{exercitadas}. Uma interface que nada implementa, um
padrão que nada instancia e um \emph{plugin} que nada carrega são, os três, a
mesma coisa --- desenho não verificado.

% =============================================================================
\subsection{Síntese}
% =============================================================================

Os quatro subprojetos deste capítulo respondem à mesma pergunta em escalas
crescentes: \emph{onde colocar a fronteira, para que a mudança seguinte custe
pouco?}

SOLID responde no nível da classe e do módulo. Os padrões respondem com soluções
nomeadas para as fronteiras que reaparecem sempre. O ECS responde recusando a
fronteira da herança e substituindo-a pela composição. E o registro de
\emph{plugins} responde levando a fronteira até o limite do processo, onde o
outro lado é literalmente um arquivo que ainda não existia quando o programa foi
compilado.

\begin{armadilha}[toda abstração deste capítulo é reversível --- e cara antes da hora]
Nenhum princípio aqui deve ser aplicado profilaticamente. OCP antes do segundo
caso é custo puro; DIP sobre uma dependência estável é indireção sem ganho; ISP
levado longe demais vira sopa de interfaces; um ECS num programa com três tipos
de entidade é complexidade sem contrapartida. A pergunta que atravessa o
capítulo inteiro --- e que os quatro subprojetos formulam com palavras
diferentes --- é sempre: \emph{esta variação já aconteceu duas vezes?}
\end{armadilha}
